\documentclass[11pt]{article}
\usepackage[margin=1in]{geometry}
\usepackage{amsmath,amssymb,hyperref}
\newcommand{\Fp}{\mathbb F_p}
\newcommand{\polylog}{\operatorname{polylog}}
\title{Two exact reductions for the absolute field unit}
\author{Project research note}
\date{18 November 2026}
\begin{document}
\maketitle

For the independently hashed slots of \texttt{target.tex}, the absolute unit
can be recovered from integer divisibility counts, with work depending on
the number of distinct exceptional magnitudes. A second reduction works
uniformly with an additional $O(b^{4/3})$ Boolean edges in the project's
local-tensor model, where $b=\lceil\log_2 p\rceil$.
Both reductions avoid the elliptic half-lists. The project's conditional
end-to-end exponent still has leading coefficient $26.2933775156$ at scale
$b/(\log_2 b)^{1/4}$.

\section{The quantity counted}
Use the odd prime-order, paired-slot setting, a finite non-two-torsion
target $T$, and the tuned dyadic window and finite offset guards in
\texttt{target.tex}. An empty slot gives unit one. Write $n_i$ for the
number of accepted abscissas in slot $i$. The terminal multiset consists
of the weights $w_T(S)=x(T)-x(S)$ away from $\pm T,O$, with the patched
values $1$ at $\pm T$ and $-1$ at $O$. Each prefix component $j$ uses the
nonzero chord scale $L(S,V)$ of that note, including the tangent and
identity branches. For each component take the product of its weights
with their oriented multiplicities. Multiply the terminal product by
the prefix products raised to $\prod_{i>j}2n_i$.

Every predicate below depends only on the component weight. Count on all
signs with the earlier affine and identity paths. On a subbox satisfying
$(2j+11)\prod_i(2\tau_i)<r$, the raw guarded tally $Z$ for a predicate
gives $\lceil Z/r\rceil$ signed occurrences. The new auxiliary assignments
are unique. Sum over boxes, then divide in the integers by two for the
terminal component and by four for a prefix component. This works for
range predicates, including ones covering almost every weight. Paid
widths use $2^{O(m\log b)}$ subboxes. Slot sizes can be counted directly
over the inverse windows in $O(\sum_i\tau_i\,\mathrm{poly}(b))$ time.
The order and the two integer moments use the existing ordinary counts.

\section{Extract the integer prime contributions directly}
Represent $w$ by an integer $1\le w<p$. Put $h=\min(w,p-w)$ and count a
minus sign when $w>p/2$. Choose a public factored positive integer $F$,
independent of the private tuple. For every $q^i$ with $1\le i\le v_q(F)$
count occurrences with $q^i\mid h$. Terms with $q^i>p/2$ may be dropped
and the corresponding exponents truncated in $F$. For each component let
$A_-$ be its negative count, $A_{q^i}$ its divisibility counts, and $N_h$
its occurrence count at magnitude $h$. The exact product is
\[
 (-1)^{A_-}\prod_{q^i}q^{A_{q^i}}
 \prod_{\substack{h:\,N_h>0\\ h\nmid F}}
       \left(\frac{h}{\gcd(h,F)}\right)^{N_h}\pmod p .
\]
Indeed the first prime product contributes precisely $\gcd(h,F)$ per
occurrence. The correction also applies when $h$ contains a higher
power of a covered prime.

Count the exceptions in $[1,\lfloor p/2\rfloor]$ by testing $F\bmod h\ne0$.
At an occupied interval query its left child and subtract from the known
parent count to obtain its right child. Traverse nonempty children depth
first and multiply each singleton correction immediately. If there are
$D$ distinct exceptional magnitudes in the component, the root and
at most $O(bD)$ further queries suffice. The interval stack has depth
$O(b)$, and counts fit in polynomially many bits. A zero root count
certifies that the displayed prime product and sign already give the
full unit. Even a dense instance can be processed without storing a
histogram, though the traversal then recovers many individual counts.

The divisibility certificates have canonical quotients. To test the
exception predicate with $F$ of bit length $L$, read its public radix
$2^b$ limbs $f_i$ from most significant to least significant. Starting
with $R_0=0$, check ordinary integer equalities
\[
 2^bR_{i-1}+f_i=Q_i h+R_i,\qquad 0\le R_i<h,\quad 0\le Q_i<2^b .
\]
All words are forced by Euclidean division. One stage uses one
$b$-by-$b$ product, carries and range checks with $O(b)$ Boolean edges
in the arithmetic compiler. Fresh copies joined on the residue limbs
keep the degree and parallel multiplicity bounds. The full predicate
adds $O(b+L)$ edges. A single $q^i\mid h$ predicate costs $O(b)$.
One can supply verified public prime powers and form $F$ using
$\mathrm{poly}(L,b)$ work. There are at most $L$ prime-power queries.
A sufficient condition to preserve the leading log-time estimate
per query is $(b+L)\log\log b=o(mb)$. This pays even for unbundled
binary links at the original simplex dimension $O(\log\log b)$.
In particular a cover with $O(b)$ bits is admissible at tuned $m$.
The exceptional searches across the components are charged separately.

\paragraph{An obstruction on the actual slots.}
There is a deterministic lower bound on the work of this particular
singleton traversal. Fix all but the largest nonempty slot, assuming
all slots are nonempty, and allow both signs in the remaining slot.
Translation gives $2n_{\max}$ distinct terminal sums. Drop the at most
three sums $\pm T,O$. Each remaining centered magnitude has at most
four preimages, two abscissas and two signs per abscissa. Consequently
\[
 D_{\rm terminal}\ \ge\
 \max\{0,\,\lceil(2n_{\max}-3)/4\rceil-\tau(F)\},
\]
where $\tau(F)$ counts positive divisors. For $F$ of $L$ bits,
$\log_2\tau(F)=O(L/\log L)$ as $L\to\infty$. To see this directly,
primes up to $\sqrt L$ contribute at most
$\sqrt L\log_2(L+1)$ to the logarithm. For larger primes the sum of
their exponents is at most $2L/\log_2 L$, and $v+1\le 2^v$.
At tuned $m=\Theta((\log b)^{1/4})$, even $L=O(mb)$ gives
$\tau(F)=2^{o(b/m)}$. Thus whenever $n_{\max}=2^{(1+o(1))b/m}$,
the number of exceptional leaves satisfies
$\log_2(1+D_{\rm terminal})\ge(1-o(1))b/m$. This last premise is on realized
slot occupancy. A width alone does not imply it. The argument concerns
enumeration of exceptions for this cover and permits other joint formulas.

\section{A uniform block certificate}
Let $g$ generate $\Fp^\times$. Write $e\in[0,p-2]$ in radix $2^s$,
with digits $u_i$ and a truncated last digit. Put
$g_i=g^{2^{is}}\bmod p$ and $k=\lceil\operatorname{bitlength}(p-2)/s\rceil$.
Certify $w=\prod_{i=0}^{k-1}h_i$ in the field, where
$h_i=g_i^{u_i}\bmod p$. There are $k-1$ modular products for $k\ge1$.
A comparison of $e$ with $p-2$ across digit blocks has forced
prefix comparison states. One can compare a block with the public
digits of $p-2$ and propagate a three-valued comparison. The range
check also applies when all blocks except the last are maximal.
Every nonzero $w$ then has exactly one exponent, with exactly one
extension to the wire variables for that exponent. Count each exponent
bit in turn. If $C_\ell$ is its oriented component count, the component
product is $g^{\sum_\ell 2^\ell C_\ell}\bmod p$.
There are $O(mb)$ kinds of pin before box splitting.
Public bases cost polynomial time by repeated squaring. A primitive
base can be supplied with an order certificate or tested by the same
canonical block count for $g^e=1$, before using it on any slots.
The latter needs only polynomially many candidates in expectation as
in \texttt{target.tex}.

\paragraph{Charging the digit copies and the field output.}
Divide each $b$-bit canonical $h_i$ into slices of at most $s$ bits.
Each slice predicate receives the $s$ or fewer digit bits and its own
output slice. It computes the indicated bits of $g_i^{u_i}\bmod p$
by ordinary powering on a local tensor entry. Recomputing the same
power at different vertices costs polynomial time per entry. No table
of size $2^s$ is assumed. Alternatively, tabulating all local powers
once costs $O(k2^s\,\mathrm{poly}(b))$ time and bits of storage. This
classical fixed-base decomposition is discussed by Brickell, Gordon,
McCurley and Wilson.\footnote{\url{https://www.microsoft.com/en-us/research/publication/fast-exponentiation-with-precomputation-algorithms-and-lower-bounds/}}

Use an equality vertex per digit bit to supply its copies to the slice
checks, range check and optional pin. Their total input wiring per
block is $O(s\lceil b/s\rceil)=O(b)$ for $1\le s\le b$.
The output slices have $b$ bits in total and can be connected bitwise
to the canonical input limbs of one fresh product module. Products
can be chained with fresh source-row owners per module and $O(b)$
bitwise links per intermediate value. A single-factor product can
link directly to the weight. Copy and subdivision relays of constant
length enforce bipartiteness and bounded parallel multiplicity, at
constant edge overhead. For example, replacing each new wire by a path
with private internal vertices permits length two between vertices of
the same color and length three between opposite colors. Equality on
the new ports has one extension. Existing arithmetic double-width
products, canonical quotient checks and row ownership are as in
\texttt{../short\_hash\_count/count.tex} and its cited compiler.

The added Boolean edge count is $O(b\lceil b/s\rceil)$.
The added degrees are bounded by
$O(s+\lceil b/s\rceil+\sqrt b\,\polylog b)$. This includes the
digit fan-out, the output-slice predicates, and the interfaces to the
arithmetic copies. Taking $s=\lceil b^{2/3}\rceil$ gives
$O(b^{4/3})$ edges and $O(b^{2/3}\polylog b)$ maximum degree.
The explicit output words themselves give $\Theta(b\lceil b/s\rceil)$
edges for this layout. Its degree therefore meets the square-root
bound for the enlarged root, up to polylogarithms. If the original
$\sqrt b\,\polylog b$ degree is imposed, $s=\lceil\sqrt b\rceil$
instead gives $O(b^{3/2})$ edges. These counts are for Boolean
edges between polynomial-time local tensor predicates. Expanding those
predicates into ordinary fan-in-two gates is a different cost model.

\paragraph{Why the smaller certificate has not improved the exponent.}
At tuned $m$ the original relation graph has $(36+o(1))mb$ leading
word bits. The $b^{4/3}$ term exceeds $mb$ by an unbounded polynomial
factor up to logarithms. It cannot be inserted into the earlier
$(18K+o(1))mb/\sqrt{\log_2 b}$ log-time estimate with its original
coefficient. Taking $s$ of order $b/m$ gives $O(mb)$ wires, but a
single unpinned digit-check vertex then sees $\Omega(b/m)$ bits.
This violates the available square-root-degree hypothesis for that
root. Enumerating the local exponent choices jointly would likewise
require charging the combinations. A new contraction bound exploiting
the block structure or a representation that combines many exceptional
magnitudes per count would change this accounting. Either route still
needs a faster full relation finder or a change in the training balance
to improve the solver bound.

\section{Checks and scope}
Run \texttt{python -B research/single\_target/verify\_block\_unit.py}
from the project root. The verifier checks bounded exponents across
five small fields, tests corrupted slice bits and also exercises a 64-bit
field and multi-limb public dividends. It compares the two
counting reductions with literal polynomial products after removing
all target zeros. Both methods are checked at two block widths or two
divisor covers respectively. Cases include empty and one-slot
inputs, infinity, tangencies, repeated target zeros, and $p=167$ with
$p-1=2\cdot83$. Both reductions return 104 and 200 for the two
$\mathbb F_{211}$ boxes with identical moments and order-30 projection.
With a full last slot on the same curve and $F=12$, the fixed-prefix
test finds 78 exceptional magnitudes against the lower bound of 49.
Exhaustive signed tuples, guarded offsets and exponent
witnesses implement the finite test oracle. The production reductions
take a guarded count callback. The checks validate the arithmetic and
rounding. The asymptotic contractor has not been benchmarked here.

\paragraph{Joint exceptional factors.}
The follow-up \texttt{exception\_product.tex} computes the correction
without visiting exceptional magnitudes. It uses one oriented count
per prime up to $\lfloor\sqrt p\rfloor$ and one sign count per component,
and gives a separate exact cancellation test. The dense 78-exception
fixture takes seven weighted counts for its terminal contribution.
\end{document}
